\documentclass[a4paper, 11pt]{article}

% PAQUETES BÁSICOS
\usepackage[utf8]{inputenc}
\usepackage[T1]{fontenc}
\usepackage[spanish]{babel}
\usepackage{geometry}
\geometry{a4paper, left=2.5cm, right=2.5cm, top=2.5cm, bottom=2.5cm}
\usepackage{times} % Usa la fuente Times New Roman
\usepackage{hyphenat} % Para mejor manejo de guiones
\usepackage{ragged2e} % Para mejor justificación de texto
\usepackage{hyperref} % Para enlaces (opcional)

% CONFIGURACIÓN DEL TÍTULO
\title{
    \Huge \textbf{Argumentario de Debate Presidencial} \\
    \large \vspace{0.5cm} Defensa de la Candidatura de Eduardo Artés Brichetti
}
\author{Electivo de Filosofía Política}
\date{\today}

% INICIO DEL DOCUMENTO
\begin{document}

\maketitle
\thispagestyle{empty} % No numera la página del título
\tableofcontents % Añade una tabla de contenido
\newpage

% --------------------------------------------------------------------
% PUNTOS CLAVE DEL PROGRAMA
% --------------------------------------------------------------------
\section{Puntos Clave del Programa de Gobierno}

% -------------------- ECONOMÍA --------------------
\subsection{Economía: Hacia un Estado Popular y Socialista}

\begin{itemize}
    \item \textbf{Rol Protagónico del Estado:} Se propone un cambio radical del modelo subsidiario actual a uno donde el Estado tenga un papel central y sea propietario de las empresas productivas más importantes del país.

    \item \textbf{Nacionalización de Sectores Estratégicos:} El programa contempla la (re)nacionalización y estatización de empresas en áreas clave como la minería (cobre, litio, hierro, silicio), la energía (generadoras eléctricas), la banca, las comunicaciones, la alimentación y el transporte.

    \item \textbf{Sistema de Propiedad Tripartito:} Se establecerá un sistema de propiedad que combinará:
    \begin{itemize}
        \item \textbf{Propiedad Estatal:} Empresas 100\% estatales en áreas estratégicas.
        \item \textbf{Propiedad Mixta:} Asociación entre el Estado y capitales privados.
        \item \textbf{Propiedad Privada:} Principalmente pequeñas y medianas empresas, protegidas de los monopolios.
        \item \textbf{Propiedad Cooperativa:} Empresas donde los trabajadores son los propietarios.
    \end{itemize}

    \item \textbf{Planificación Centralizada:} La economía se organizará a través de sistemas quinquenales de planificación nacional, buscando la industrialización y el desarrollo soberano del país.

    \item \textbf{Reforma Tributaria Progresiva:} Se implementará un nuevo sistema tributario que aumentará la carga impositiva sobre las grandes fortunas, las herencias y las grandes empresas, mientras se alivia a los trabajadores y pequeños comerciantes.

    \item \textbf{Seguridad Alimentaria y Energética:} Se busca la soberanía en estos ámbitos a través de la estatización de agroindustrias y grandes generadoras de energía, incluyendo el desarrollo de energía nuclear.

    \item \textbf{Política Exterior Antiimperialista:} Se buscará una política exterior soberana, fortaleciendo lazos con países que avanzan en la construcción del socialismo y se oponen al imperialismo.

    \item \textbf{Asamblea Constituyente Genuina:} Se convocará a una asamblea constituyente con participación directa de organizaciones sociales para redactar una nueva Constitución "en una hoja en blanco".
\end{itemize}

% -------------------- SEGURIDAD --------------------
\subsection{Seguridad: Nuevo Orden Institucional y Defensa de la Soberanía}
\begin{itemize}
    \item \textbf{Refundación de las Fuerzas Armadas y de Orden:} Se propone una reestructuración completa de las fuerzas de seguridad, con un carácter popular, patriótico y antiimperialista. Se introducirá un escalafón único y se les dará participación política en la Asamblea Popular.

    \item \textbf{Soberanía y Defensa Nacional:} La defensa nacional será un pilar fundamental, orientada a la independencia económica y a la protección contra amenazas imperialistas.

    \item \textbf{Combate al Crimen Organizado:} Se enfrentará con decisión al crimen organizado, considerado un enemigo de la clase trabajadora. El programa propone la pena de muerte para grandes narcotraficantes y en casos de pedofilia.

    \item \textbf{Control de la Inmigración:} Se establecerá un control estricto de la inmigración para combatir el crimen organizado internacional y proteger la soberanía nacional.

    \item \textbf{Sistema Penitenciario Centrado en la Reinserción:} Se busca sustituir el actual sistema punitivo por uno enfocado en la educación, reformación y preparación del presidiario para su reinserción en la sociedad.

    \item \textbf{Justicia Popular:} Se propone la destitución de la Corte Suprema y su reemplazo por un "Tribunal Democrático del País", con miembros designados por la "Asamblea de los Pueblos".
\end{itemize}

% -------------------- EDUCACIÓN --------------------
\subsection{Educación: Formación para una Nueva Sociedad}
\begin{itemize}
    \item \textbf{Educación Pública, Laica y Gratuita:} El Estado garantizará el derecho a una educación de calidad en todos sus niveles, desde la primaria hasta la superior.

    \item \textbf{Orientación Científica y Humanista:} El sistema educativo impartirá una visión científica y humanista del mundo, vinculada al proceso de industrialización del país.

    \item \textbf{Fin de la Mercantilización:} Se terminará con la mercantilización de la investigación universitaria, que actualmente beneficia a grandes monopolios.

    \item \textbf{Plan Nacional de Educación:} Se diseñará un nuevo currículum y estructura educativa en todos los niveles, en consonancia con las necesidades del proyecto de país.

    \item \textbf{Fortalecimiento de la Educación Técnico-Profesional:} Se potenciarán los colegios técnicos para integrarlos directamente al sector productivo nacional.

    \item \textbf{Participación de las Comunidades Educativas:} Se dará un rol central a los gremios de profesores, sindicatos, organizaciones de trabajadores de la educación y estudiantes organizados en la formulación de los principios de la educación.
\end{itemize}
\newpage

% --------------------------------------------------------------------
% PREGUNTAS PARA OTROS CANDIDATOS
% --------------------------------------------------------------------
\section{Preguntas para Representantes de Otros Candidatos}

% -------------------- KAISER --------------------
\subsection{Para el representante de Johannes Kaiser (Ideas Libertarias)}
\begin{itemize}
    \item \textbf{Sobre Temas de Género y Violencia:} Johannes Kaiser tiene un historial de declaraciones misóginas y violentas, como sugerir que "las feminazis necesitan un buen pico militar" o que a las mujeres se les puede "hacer evangélicas a puros pichulazos". ¿Sigue sosteniendo estas ideas que promueven la violencia sexual como herramienta política? ¿Cómo puede un candidato con estas posturas garantizar los derechos y la seguridad de más de la mitad de la población de Chile, que son mujeres?

    \item \textbf{Sobre el Rol del Estado:} Ustedes proponen una reducción drástica del Estado. ¿Cómo pretenden garantizar el bienestar y la seguridad de la población con un Estado mínimo, especialmente en salud, educación y pensiones, que históricamente han requerido inversión pública para asegurar un acceso equitativo?
    
    \item \textbf{Sobre la Desigualdad:} Su enfoque se centra en la libertad de mercado. El programa de Artés considera la desigualdad una consecuencia de la explotación capitalista. ¿No creen que la extrema concentración de la riqueza, como la de Chile, socava la misma libertad que ustedes defienden al limitar las oportunidades reales para la mayoría?
\end{itemize}

% -------------------- KAST --------------------
\subsection{Para el representante de José Antonio Kast (Derecha Conservadora)}
\begin{itemize}
    \item \textbf{Sobre Gasto Público y "Parásitos":} Su programa promete un recorte de 6 mil millones de dólares en gasto público, eliminando lo que un asesor clave de su campaña ha denominado "parásitos del Estado", sin supuestamente tocar el gasto social. ¿Puede especificar qué programas y qué tipo de funcionarios públicos serían eliminados? ¿Y cómo garantiza que este recorte masivo no afectará directamente la capacidad del Estado para responder a las necesidades de la ciudadanía, más allá de la retórica?

    \item \textbf{Sobre Soberanía y Recursos Naturales:} Su programa defiende la inversión extranjera. El de Artés propone renacionalizar nuestros recursos estratégicos como el cobre y el litio. ¿Por qué deberíamos seguir permitiendo que las transnacionales se lleven nuestras riquezas, en lugar de usar esos recursos para construir un país más justo para todos los chilenos?

    \item \textbf{Sobre Orden Público y Derechos Humanos:} Usted habla de "mano dura". ¿Cómo nos garantiza que su enfoque en el orden no terminará, como ha ocurrido históricamente, en una represión indiscriminada contra los movimientos sociales que luchan por sus derechos, criminalizando la protesta?
\end{itemize}

% -------------------- MATTHEI --------------------
\subsection{Para el representante de Evelyn Matthei (Derecha Tradicional)}
\begin{itemize}
    \item \textbf{Sobre Crecimiento Económico:} Su programa se centra en el crecimiento a través de incentivos a la inversión privada. ¿De qué le sirve a la mayoría un crecimiento que se concentra en pocas manos y no se traduce en mejores salarios, pensiones dignas y servicios públicos de calidad?

    \item \textbf{Sobre la Modernización del Estado:} Usted propone modernizar el Estado, pero manteniendo su rol subsidiario. ¿No es la idea de un "Estado eficiente" pero subsidiario una contradicción? ¿Cómo puede ser eficiente en resolver los grandes problemas del país si se le prohíbe participar activamente en la economía?

    \item \textbf{Sobre Educación:} Usted habla de fortalecer la educación técnica dentro del actual sistema de mercado. ¿No es su propuesta una forma de mantener la educación como un negocio y un filtro social, en lugar de transformarla en una verdadera herramienta de liberación y desarrollo para todo el pueblo de Chile?
\end{itemize}

\end{document}
